% Importado desde: 2026-03-28_QW_3D_Weyl_a_Dirac_3p1_consistente.tex
\chapter{Derivacion Pedagogica: --- de dos bloques Weyl discretos a una construccion --- consistente de Dirac en \(3+1\)}
\label{ch:qw-3d-weyl-a-dirac-3p1}

\section{Objetivo}
En el capitulo anterior construimos una \textit{quantum walk} minima en \(3D\) cuyo limite infrarrojo produce un bloque de Weyl isotropo:
\begin{equation}
H_W = v(\sigma_x p_x+\sigma_y p_y+\sigma_z p_z).
\end{equation}

El siguiente paso solo tiene sentido si es \textbf{algebraicamente consistente}. Es decir, no basta decir \enquote{ponemos dos Weyl juntos y agregamos una masa}; hay que verificar que:
\begin{enumerate}[label=\arabic*.]
    \item la dinamica discreta ampliada puede escribirse de forma natural;
    \item el limite efectivo produce exactamente un Hamiltoniano de Dirac de cuatro componentes;
    \item la masa anticommuta con el bloque cinetico;
    \item y de ahi emerge una realizacion explicita de la algebra de Clifford.
\end{enumerate}

Ese es el objetivo de este capitulo.

\section{Primer bloque: recordar el sector Weyl}
Partimos de la caminata discreta elemental
\begin{equation}
U_W = U_z U_y U_x,
\qquad
U_i = \ee^{-\,\ii a p_i \sigma_i}.
\label{eq:UW}
\end{equation}

Para \(a|\vp|\ll 1\), esta produce
\begin{equation}
U_W \simeq \bbI - \ii \Delta t\,H_W,
\qquad
H_W = v(\sigma_x p_x+\sigma_y p_y+\sigma_z p_z),
\label{eq:HW}
\end{equation}
con
\begin{equation}
v=\frac{a}{\Delta t}.
\end{equation}

Este es el bloque de quiralidad positiva.

\section{Segundo bloque: construir la quiralidad opuesta}
Para pasar a Dirac necesitamos otro bloque con orientacion opuesta. La forma mas limpia es
\begin{equation}
H_{-W} = -\,v(\sigma_x p_x+\sigma_y p_y+\sigma_z p_z).
\label{eq:Hminus}
\end{equation}

\subsection{Version unitaria}
La caminata correspondiente puede escribirse como
\begin{equation}
U_{-W} = U_x^{-1} U_y^{-1} U_z^{-1},
\label{eq:UmW}
\end{equation}
porque, a primer orden,
\begin{equation}
U_i^{-1} = \ee^{+\,\ii a p_i \sigma_i}
\simeq \bbI + \ii a p_i \sigma_i.
\end{equation}

Entonces
\begin{equation}
U_{-W}\simeq \bbI + \ii a(\sigma_x p_x+\sigma_y p_y+\sigma_z p_z)
= \bbI - \ii \Delta t\,H_{-W}.
\end{equation}

\subsection{Leccion}
No estamos inventando el segundo bloque a mano: sale de una dinamica discreta igualmente local, pero con orientacion opuesta.

\begin{figure}[h!]
\centering
\begin{tikzpicture}[>=Latex,node distance=1.2cm]
  \node[draw,rounded corners,fill=blue!7,minimum width=3.2cm,minimum height=0.95cm,align=center] (a) {bloque Weyl \(+\)\\\(U_zU_yU_x\)};
  \node[draw,rounded corners,fill=blue!7,minimum width=3.2cm,minimum height=0.95cm,align=center,right=4cm of a] (b) {bloque Weyl \(-\)\\\(U_x^{-1}U_y^{-1}U_z^{-1}\)};
  \draw[<->,thick] (a) -- node[above] {\small quiralidad opuesta} (b);
\end{tikzpicture}
\caption{Los dos bloques discretos que usaremos para construir el sector Dirac.}
\end{figure}

\section{Tercer paso: ensamblar el espacio de cuatro componentes}
Ahora introducimos un nuevo espacio interno de dos componentes, descrito por matrices \(\tau_i\), para distinguir ambos sectores Weyl.

El espinor ampliado es
\begin{equation}
\Psi=
\begin{pmatrix}
\psi_+\\
\psi_-
\end{pmatrix},
\end{equation}
donde \(\psi_+\) y \(\psi_-\) son espinores de dos componentes.

\subsection{Operador unitario sin masa}
La evolucion conjunta sin masa puede escribirse, a primer orden, como
\begin{equation}
U_0 \simeq \bbI - \ii \Delta t\,H_0,
\end{equation}
con
\begin{equation}
H_0=
\begin{pmatrix}
H_W & 0\\
0 & H_{-W}
\end{pmatrix}.
\label{eq:H0block}
\end{equation}

Usando \eqref{eq:HW} y \eqref{eq:Hminus}, esto equivale a
\begin{equation}
H_0 = v\,\tau_z\otimes(\sigma_x p_x+\sigma_y p_y+\sigma_z p_z).
\label{eq:H0tau}
\end{equation}

\subsection{Interpretacion}
El signo de \(\tau_z\) separa exactamente los dos sectores de quiralidad opuesta.

\section{Cuarto paso: introducir la masa como mezcla local}
La masa debe mezclar ambos bloques. La forma minima es
\begin{equation}
H_m = m\,\tau_x\otimes \bbI_2.
\label{eq:Hm}
\end{equation}

\subsection{Version discreta}
En lenguaje de caminata, esto corresponde a una moneda interna adicional
\begin{equation}
M = \ee^{-\,\ii \Delta t\,m\,\tau_x\otimes \bbI_2}.
\label{eq:M}
\end{equation}

Tomando primero el paso sin masa y luego la mezcla:
\begin{equation}
U_D = M\,U_0.
\label{eq:UD}
\end{equation}

Para \(\Delta t\) peque\~no:
\begin{equation}
M \simeq \bbI - \ii \Delta t\,m\,\tau_x\otimes\bbI_2.
\end{equation}

Por tanto,
\begin{equation}
U_D \simeq \bbI - \ii \Delta t(H_0+H_m),
\end{equation}
y el Hamiltoniano efectivo total es
\begin{equation}
H_D = v\,\tau_z\otimes(\sigma_x p_x+\sigma_y p_y+\sigma_z p_z)
+ m\,\tau_x\otimes \bbI_2.
\label{eq:HD}
\end{equation}

\begin{figure}[h!]
\centering
\begin{tikzpicture}[>=Latex,node distance=1.1cm]
  \node[draw,rounded corners,fill=blue!7,minimum width=2.8cm,minimum height=0.88cm,align=center] (w1) {Weyl \(+\)};
  \node[draw,rounded corners,fill=blue!7,minimum width=2.8cm,minimum height=0.88cm,align=center,right=3.2cm of w1] (w2) {Weyl \(-\)};
  \node[draw,rounded corners,fill=orange!10,minimum width=2.8cm,minimum height=0.88cm,align=center,below=1.45cm of $(w1)!0.5!(w2)$] (m) {mezcla \(m\,\tau_x\)};
  \node[draw,rounded corners,fill=green!8,minimum width=3.5cm,minimum height=0.95cm,align=center,below=1.45cm of m] (d) {Dirac \(3+1\)};
  \draw[->,thick] (w1) -- (d);
  \draw[->,thick] (w2) -- (d);
  \draw[->,thick] (m) -- (d);
\end{tikzpicture}
\caption{La masa no es un adorno: es el acoplamiento local que convierte dos bloques Weyl en un sector Dirac.}
\end{figure}

\section{Verificacion de consistencia algebraica}
Ahora viene la parte decisiva.

\subsection{Definiciones}
Definimos
\begin{equation}
\alpha_i = \tau_z\otimes \sigma_i,
\qquad
\beta = \tau_x\otimes \bbI_2.
\label{eq:alphabeta}
\end{equation}

Entonces \eqref{eq:HD} toma la forma estandar
\begin{equation}
H_D = v\sum_{i=1}^3 \alpha_i p_i + m\beta.
\label{eq:HDstandard}
\end{equation}

\subsection{Cuadrados}
Como \(\tau_z^2=\tau_x^2=\bbI_2\) y \(\sigma_i^2=\bbI_2\), tenemos
\begin{equation}
\alpha_i^2=\bbI_4,
\qquad
\beta^2=\bbI_4.
\end{equation}

\subsection{Anticonmutacion espacial}
Para \(i\neq j\),
\begin{equation}
\{\alpha_i,\alpha_j\}
=
\tau_z^2\otimes\{\sigma_i,\sigma_j\}
=0,
\end{equation}
y para \(i=j\),
\begin{equation}
\{\alpha_i,\alpha_i\}=2\bbI_4.
\end{equation}

Luego
\begin{equation}
\{\alpha_i,\alpha_j\}=2\delta_{ij}\bbI_4.
\label{eq:aa}
\end{equation}

\subsection{Anticonmutacion con la masa}
Usamos \(\{\tau_z,\tau_x\}=0\):
\begin{equation}
\{\alpha_i,\beta\}
=
\{\tau_z\otimes\sigma_i,\tau_x\otimes\bbI_2\}
=
\{\tau_z,\tau_x\}\otimes \sigma_i
=0.
\label{eq:ab}
\end{equation}

\subsection{Conclusion algebraica}
Las relaciones \eqref{eq:aa} y \eqref{eq:ab} son exactamente las del Hamiltoniano de Dirac. Por tanto, la construccion es consistente.

\begin{displayquote}
\textbf{la masa propuesta no destruye el bloque Weyl; lo completa en una realizacion de cuatro componentes que satisface la algebra correcta.}
\end{displayquote}

\section{Chequeo del espectro}
La mejor prueba de consistencia es cuadrar el Hamiltoniano:
\begin{align}
H_D^2
&=
v^2\sum_{i,j}\alpha_i\alpha_j p_i p_j
+ mv\sum_i(\alpha_i\beta+\beta\alpha_i)p_i
+ m^2\beta^2
\notag\\
&=
v^2(p_x^2+p_y^2+p_z^2)\bbI_4 + m^2\bbI_4,
\end{align}
porque los terminos cruzados desaparecen por anticonmutacion.

Luego:
\begin{equation}
E(\vp)=\pm\sqrt{v^2(p_x^2+p_y^2+p_z^2)+m^2}.
\label{eq:DiracSpec}
\end{equation}

Esto coincide exactamente con el espectro de Dirac esperado.

\section{Lectura en terminos de matrices gamma}
Una vez identificado \(\alpha_i\) y \(\beta\), definimos
\begin{equation}
\gamma^0=\beta,
\qquad
\gamma^i=\beta\alpha_i.
\label{eq:gamma}
\end{equation}

Con nuestras matrices:
\begin{equation}
\gamma^0=\tau_x\otimes\bbI_2,
\qquad
\gamma^i=-\,\ii \tau_y\otimes \sigma_i.
\end{equation}

Estas satisfacen
\begin{equation}
\{\gamma^\mu,\gamma^\nu\}=2\eta^{\mu\nu}\bbI_4,
\qquad
\eta^{\mu\nu}=\mathrm{diag}(1,-1,-1,-1).
\end{equation}

O sea: la construccion no solo es consistente a nivel hamiltoniano, sino que ya instala una realizacion explicita de la estructura de Clifford \(Cl_{1,3}\).

\section{Que simetrias conserva y cuales no}
\subsection{Lo que si conserva}
La construccion mantiene:
\begin{enumerate}[label=\arabic*.]
    \item localidad del paso;
    \item estructura discreta de traslaciones;
    \item el uso de bloques Weyl con orientacion controlada;
    \item y la posibilidad de recuperar isotropia efectiva en el infrarrojo.
\end{enumerate}

\subsection{Lo que aun no garantiza}
No garantiza todavia, por si sola:
\begin{enumerate}[label=\arabic*.]
    \item la emergencia exacta de todo el grupo de Lorentz;
    \item el control global de dobladores en toda la zona de Brillouin;
    \item ni la unicidad de esta construccion como realizacion microscópica.
\end{enumerate}

Eso esta bien: en esta etapa solo queriamos consistencia del paso \enquote{Weyl discreto \(\to\) Dirac \(3+1\)}.

\section{Por que este paso si es coherente con el programa}
Ahora ya podemos responder a tu pregunta implicita: \textbf{si, este paso es consistente con la ruta del proyecto}.

\subsection{Por que}
Porque respeta exactamente el orden que nos habiamos impuesto:
\begin{enumerate}[label=\arabic*.]
    \item primero un bloque discreto tipo Weyl;
    \item luego su quiralidad;
    \item luego el doblamiento controlado en un espacio interno mayor;
    \item y solo entonces la masa y la estructura de Clifford.
\end{enumerate}

No hay aqui ningun salto arbitrario ni una identificacion apresurada. El paso a Dirac aparece como prolongacion natural del bloque Weyl.

\begin{figure}[h!]
\centering
\begin{tikzpicture}[
    box/.style={draw, rounded corners, thick, align=center, minimum width=3.7cm, minimum height=1.02cm},
    arr/.style={-{Latex[length=3mm]}, thick}
]
\node[box, fill=blue!7] (a) at (0,0) {QW Weyl \(+\)};
\node[box, fill=blue!7] (b) at (4.5,0) {QW Weyl \(-\)};
\node[box, fill=yellow!12] (c) at (9,0) {mezcla local\\de masa};
\node[box, fill=green!8] (d) at (13.5,0) {Dirac \(3+1\)\\consistente};
\draw[arr] (a) -- (c);
\draw[arr] (b) -- (c);
\draw[arr] (c) -- (d);
\end{tikzpicture}
\caption{La consistencia del paso Weyl \(\to\) Dirac queda asegurada por la estructura de bloques y la anticonmutacion del termino de masa.}
\end{figure}

\section{Mapa del siguiente paso}
Con esta pieza cerrada, el siguiente paso natural ya no es preguntar \enquote{si se puede}. Eso ya quedo mostrado a nivel efectivo.

La pregunta que sigue es mas fina:
\begin{displayquote}
\textbf{como se insertan Nielsen--Ninomiya, la estructura global de la zona de Brillouin y las simetrias discretas completas dentro de esta construccion de cuatro componentes?}
\end{displayquote}

En otras palabras: el siguiente bloque serio deberia estudiar la \textbf{consistencia global}, no solo local.

\section{Resumen final}
Tomando dos bloques discretos de Weyl con quiralidad opuesta y una mezcla local de masa, construimos una dinamica ampliada de cuatro componentes cuyo Hamiltoniano efectivo es
\begin{equation}
H_D = v\,\tau_z\otimes(\vsigma\cdot\vp) + m\,\tau_x\otimes\bbI_2.
\end{equation}

Verificamos paso a paso que:
\begin{enumerate}[label=\arabic*.]
    \item la construccion viene de un paso unitario discreto natural;
    \item la masa anticommuta con el bloque cinetico;
    \item el espectro es el de Dirac;
    \item y la algebra de Clifford \(Cl_{1,3}\) aparece de forma explicita.
\end{enumerate}

Eso hace que este sea, efectivamente, un paso consistente dentro del programa. No cierra toda la investigacion, pero si cierra de forma limpia el puente local entre \textit{Weyl discreto} y \textit{Dirac \(3+1\)}.
